\documentclass[11pt]{article}
\usepackage[T1]{fontenc}
\usepackage[utf8]{inputenc}
\usepackage{lmodern,amsmath,amssymb,amsthm,booktabs,array,microtype}
\usepackage[margin=1in]{geometry}
\usepackage[colorlinks=true,linkcolor=blue,citecolor=blue,urlcolor=blue]{hyperref}
\usepackage[numbers,sort&compress]{natbib}
\hypersetup{pdftitle={Boundary log-concavity of binary reciprocal-subsum numerators},
 pdfauthor={Xiaojian Zeng}}
\newtheorem{theorem}{Theorem}
\newtheorem{lemma}[theorem]{Lemma}
\newtheorem{proposition}[theorem]{Proposition}
\newcommand{\B}{\mathcal B}
\newcommand{\T}{\mathcal T}

\newcommand{\ind}{\mathbf1}
\DeclareMathOperator{\degp}{deg}
\title{Boundary log-concavity of binary reciprocal-subsum numerators}
\author{Xiaojian Zeng\\[3pt]
\small PBC School of Finance, Tsinghua University, 100190\\
\small\texttt{xiaojianzeng@foxmail.com}}
\date{}
\begin{document}
\maketitle
\begin{abstract}
Let $N_n(x)$ be the binary reciprocal-subsum numerator introduced by
Ballantine, Beck, Feigon and Maurischat. We prove strict coefficient
log-concavity, uniformly for every fixed $0<\varepsilon<1$ and sufficiently
large $n$, at the first
$\lfloor(1-\varepsilon)\sqrt{\log_2 n/2}\rfloor$ positions and their
reflections. An elementary bound for the mean multiplicity of ones and two
coefficient envelopes also give an explicit rational sufficient test without
enumerating binary partitions. For every $n\ge6$, the first three positions
and their reflections satisfy strict inequalities; at even $n$ we prove a
stronger comparison with the first recurrence summand. The third-position
proof uses exact polynomial identities on a complete set of parameter
chambers. The full log-concavity conjecture remains unresolved here.
\end{abstract}

\section{The polynomial family and the boundary results}
Let $\B(n)$ be the set of partitions of $n$ into powers of two, and let
$m_\lambda(q)$ denote the multiplicity of $q$ in $\lambda$.
Write $b(n)=|\B(n)|$, including $b(0)=1$. For $n\ge1$, put
\begin{equation}\label{eq:def}
 N_n(x)=\sum_{\lambda\in\B(n)}h_{n,\lambda}(x),\qquad
 h_{n,\lambda}(x)=\prod_{2^k\le n}
 (1+x^{2^k})^{\lfloor n/2^k\rfloor-m_\lambda(2^k)},
\end{equation}
and set $N_0=1$. This is precisely $\operatorname{num}_{\B}(n,x)$ of
\cite[Remark 6]{BBFM}: the common divisor of the individual summands used
in that normalization is $1$.
Let $c_{n,j}=[x^j]N_n$ and
\[
 d(n)=\degp N_n=\sum_{k\ge1}2^k\lfloor n/2^k\rfloor.
\]
For $P(x)=\sum p_jx^j$, define
$\T_j(P)=p_j^2-p_{j-1}p_{j+1}$.
An internal position means $1\le j<\degp P$; no parity subsequence or
coefficient normalization is used.

The originating paper proves positivity of all coefficients,
palindromicity, and $N_{2m+1}=N_{2m}$
\cite[Proposition 4.1]{BBFM}. Each summand in \eqref{eq:def} has degree
$d(n)$ and is palindromic, which also explains the reflection symmetry.
The conjecture under consideration is
\begin{equation}\label{eq:original}
 \T_j(N_n)\ge0\qquad(n\ge6,\ 1\le j<d(n)),
\end{equation}
as stated in \cite[Conjecture 7]{BBFM}.
The remark in Section 1.2.2 of \cite[version 2]{COZ} records that the
failures at $n=4,5$ concern an earlier range and leave the corrected question
unresolved. Our results are restricted to boundary positions.

For even $n$, write
\begin{equation}\label{eq:rec}
 N_n=A_n+B_n,\qquad A_n=(1+x)^nN_{n/2}(x^2),\qquad
 B_n=\prod_{i=1}^{v_2(n)}(1+x^{2^i})N_{n-2}(x),
\end{equation}
where $v_2(n)$ is the exponent of $2$ in $n$.
This is \cite[Theorem 4.7]{BBFM}. Both summands have degree $d(n)$.
The recurrence separates partitions with no ones, whose parts can be
halved, from those with at least two ones, from which two ones can be deleted.

\begin{theorem}\label{thm:main}
For the numerators \eqref{eq:def}, the following assertions hold.
\begin{enumerate}
\item For every real $0<\varepsilon<1$ there is an integer $n_\varepsilon$
such that, for every integer $n\ge n_\varepsilon$,
\[
 \T_j(N_n)>0\quad\text{for }1\le j\le
 \left\lfloor(1-\varepsilon)\sqrt{\frac{\log_2 n}{2}}\right\rfloor
 \quad\text{and at each }d(n)-j.
\]
\item For every integer $n\ge6$, strict positivity holds at
$j=1,2,3,d(n)-1,d(n)-2,d(n)-3$. For even $n\ge6$, at these positions,
\[
 \T_j(N_n)>\T_j(A_n)>0.
\]
\item Define the rational number
\begin{equation}\label{eq:M}
 M_n=\min_{\substack{K\in\mathbb Z\\0\le K\le\lfloor\log_2n\rfloor}}
 \frac{n+2^{K+1}-K-2}{K+1}.
\end{equation}
For $n\ge2$ and an internal integer position $1\le j\le n$, the conditions
$jM_n<n$ and
\begin{equation}\label{eq:test}
 \binom nj^2\left(1-\frac{jM_n}{n}\right)^2>
 \binom{n+j-2}{j-1}\binom{n+j}{j+1}
\end{equation}
imply $\T_j(N_n)>0$.
\end{enumerate}
\end{theorem}

The added conclusion relative to the structural and initial-coefficient
formulas in \cite[Section 4]{BBFM} is a uniform range whose width tends
to infinity, together with a count-free sufficient test and the three-position
comparison for all admissible $n$. The growing range requires control of
coefficient errors while $j$ grows; a fixed-index limit does not suffice.
No equivalent boundary statement was found in the primary sources inspected
for this note as of 13 September 2026. This is a bounded literature assessment,
not a claim of exhaustive priority. Neither the width constant nor the
sufficient criterion is asserted to be optimal.

\section{Occupation estimates and the growing range}
Choose $\lambda\in\B(n)$ uniformly. For a power of two $q$, let
$\mu_q=\mathbb E m_\lambda(q)$ and put $\mu=\mu_1$.

\begin{lemma}\label{lem:mean}
For $n\ge2$, one has $\mu\le M_n$. If $L=\log_2n$ and $n\ge4$, then
\begin{equation}\label{eq:mean}
 \frac{\mu}{n}\le\frac{M_n}{n}
 \le\frac{1+2/L}{L-\log_2L}=\frac{1+o(1)}L.
\end{equation}
\end{lemma}
\begin{proof}
Adding a one shows that $b(h)$ is nondecreasing for $h\ge0$; extend it by
$b(h)=0$ for $h<0$. Deleting $r$ copies of $q$ gives
\[
 \mu_q=\sum_{r\ge1}\frac{b(n-rq)}{b(n)}.
\]
In the corresponding sum for $\mu$, isolate its first $q-1$ terms, each
at most $1$, and group the rest into blocks $rq,\ldots,(r+1)q-1$.
Monotonicity bounds each block by $qb(n-rq)/b(n)$. Hence
$\mu\le q-1+q\mu_q$.
Since the total partition weight has expectation $n$, for every admissible $K$,
\[
 n=\sum_{k\ge0}2^k\mu_{2^k}
 \ge (K+1)\mu-\sum_{k=0}^K(2^k-1)
 =(K+1)\mu-(2^{K+1}-K-2).
\]
Minimizing proves the first assertion.
For the second, take $K=\lfloor L-\log_2L\rfloor$.
Then $2^{K+1}\le2n/L$ and $K+1\ge L-\log_2L>0$.
Dropping the nonpositive term $-K-2$ in the numerator gives \eqref{eq:mean}.
\end{proof}

\begin{lemma}\label{lem:envelopes}
For $0\le j\le n$,
\begin{align}
 c_{n,j}&\ge b(n)\binom nj(1-j\mu/n),\label{eq:lower1}\\
 c_{n,j}&\ge b(n)\frac{n^j}{j!}(1-\delta_{n,j}),\qquad
 \delta_{n,j}=\frac{j\mu}{n}+\frac{j(j-1)}{2n}.\label{eq:lower2}
\end{align}
For every $j\ge0$ and $n\ge1$,
\begin{equation}\label{eq:upper}
 c_{n,j}\le b(n)\binom{n+j-1}{j}
 \le b(n)\frac{n^j}{j!}\exp\left(\frac{j(j-1)}{2n}\right).
\end{equation}
Here coefficients outside the polynomial's degree range are zero.
\end{lemma}
\begin{proof}
The factor $(1+x)^{n-m_\lambda(1)}$ in each summand gives the lower bound
$\binom{n-m_\lambda(1)}j$. A uniform $j$-subset of $n$ labels avoids a
specified set of $m$ forbidden labels with probability at least $1-jm/n$,
by the union bound. Averaging over $m=m_\lambda(1)$ proves
\eqref{eq:lower1}. Alternatively draw $j$ labels independently. The
probability of a forbidden draw or a collision is at most
$jm/n+\binom j2/n$. The number of ordered acceptable tuples is
$j!\binom{n-m}j$, which proves \eqref{eq:lower2} after averaging.
These lower bounds remain valid when their right sides are negative.

Let $\preceq$ denote coefficientwise order of formal series. Every exponent
in \eqref{eq:def} is between $0$ and $n$. Unique finite binary expansion gives
\[
 h_{n,\lambda}(x)\preceq
 \prod_{k\ge0}(1+x^{2^k})^n=(1-x)^{-n}.
\]
Only finitely many factors contribute to a fixed coefficient, so this uses
no analytic infinite-product limit. Summing over $\lambda$ proves the first
upper bound. The second follows from
$\prod_{r=0}^{j-1}(1+r/n)\le\exp(\sum_{r=0}^{j-1}r/n)$.
\end{proof}

\begin{proof}[Proof of Theorem \ref{thm:main}(3)]
The condition $jM_n<n$ makes the lower envelope in \eqref{eq:lower1}
positive. Square it, replace $\mu$ by its upper bound $M_n$, and apply the
first bound in \eqref{eq:upper} at $j-1$ and $j+1$.
After canceling $b(n)^2$, \eqref{eq:test} gives the claimed strict inequality.
\end{proof}

\begin{proof}[Proof of Theorem \ref{thm:main}(1)]
Fix $\alpha=1-\varepsilon$ and set
$J_n=\lfloor\alpha\sqrt{L/2}\rfloor$, where $L=\log_2n$.
Since $d(n)\ge2\lfloor n/2\rfloor\ge n-1$, these positions are internal
and $J_n<n$ for sufficiently large $n$.
Uniformly for $1\le j\le J_n$, Lemma \ref{lem:mean} yields
\begin{equation}\label{eq:uniform}
 \delta_{n,j}=o(1),\quad
 j\delta_{n,j}\le\alpha^2/2+o(1),\quad
 j\delta_{n,j}^2=o(1),\quad
 \frac{j(j^2-j+1)}n=o(1).
\end{equation}
Indeed, $j^2\mu/n\le\alpha^2/2+o(1)$ and the remaining contribution to
$j\delta_{n,j}$ is $O(L^{3/2}/n)$. Also
$\delta_{n,j}=O(L^{-1/2})$, proving the other assertions.
The positive lower envelope and the upper envelopes give
\begin{equation}\label{eq:ratio}
 \frac{c_{n,j}^2}{c_{n,j-1}c_{n,j+1}}
 \ge R_{n,j}:=\frac{j+1}{j}(1-\delta_{n,j})^2
 \exp\left(-\frac{j^2-j+1}{n}\right).
\end{equation}
Using
$\log(1+x)\ge x-x^2/2$ for $x\ge0$ and
$\log(1-y)\ge-y-y^2/(2(1-y))$ for $0\le y<1$, equations
\eqref{eq:uniform}--\eqref{eq:ratio} imply, with a uniform remainder,
\[
 j\log R_{n,j}\ge1-\frac1{2j}-\alpha^2-o(1).
\]
Choose a fixed integer $J$ with $1/(2J)<(1-\alpha^2)/2$.
For all sufficiently large $n$ this is positive simultaneously for
$J\le j\le J_n$. For the finitely many $1\le j<J$, one has
$R_{n,j}\to(j+1)/j>1$. Enlarging the same threshold handles these as well.
Palindromicity gives the reflected inequalities.
\end{proof}

\section{The first three positions for every admissible parameter}
Put $n=2m$, $w_k=b(k)$, $w=w_m$, and
\[
 S_r=\sum_{k=0}^m k^rw_k,\qquad R_r=\sum_{k=0}^{m-1}k^rw_k,
\]
with $k^0=1$ at $k=0$. Write $c_j=[x^j]N_{2m}$ and $a_j=[x^j]A_{2m}$.
The coefficient identities needed below are
\begin{align}
 c_0&=S_0,&c_1&=2S_1,&c_2&=2S_2,\label{eq:low}\\
 c_3&=\sum_{k=0}^m H_m(k)w_k,&
 H_m(k)&=\frac{4k^3+2k}{3}-(m-k)(m-k+1).\label{eq:cubic}
\end{align}
The first three formulas are prior identities
\cite[Corollary 4.5 and Remark 8]{BBFM}.
To obtain the next one, first note that deletion of ones gives
\begin{equation}\label{eq:linear}
 [x]N_k=kw_k-\sum_{h<k}w_h,
\end{equation}
also \cite[Corollary 4.2]{BBFM}. The carry factor in \eqref{eq:rec} is
$1+x^2+O(x^4)$. Its cubic increment is therefore
\[
 \binom{2k}{3}w_k+2k[x]N_k+[x]N_{2k-2}
 =\frac{4k^3+2k}{3}w_k-2\sum_{h<k}(k-h)w_h.
\]
Summing over $k=1,\ldots,m$ gives \eqref{eq:cubic}.
Direct extraction in $A_{2m}$ gives
\begin{align*}
 a_0&=w,&a_1&=2mw,&a_2&=2m^2w-R_0,&
 a_3&=\frac{4m^3+2m}{3}w-2mR_0.
\end{align*}
Consequently
\begin{equation}\label{eq:apos}
 \T_1(A_{2m})=2m^2w^2+wR_0>0,\qquad
 \T_2(A_{2m})=\frac{4m^2(m^2-1)}3w^2+R_0^2>0
\end{equation}
for $m\ge3$.

\subsection{First and second positions}
Expanding \eqref{eq:low} gives
\begin{align}
 \T_1(N_{2m})-\T_1(A_{2m})
 &=w\sum_{k<m}g_m(k)w_k+4R_1^2-2R_0R_2,\label{eq:first}\\
 g_m(k)&=8mk-2k^2-2m^2-1.\notag
\end{align}
For increasing sequences $f_k,w_k$, summing
$(f_i-f_j)(w_i-w_j)\ge0$ gives Chebyshev's sum inequality.
It implies $2R_1\ge(m-1)R_0$. Since $R_2\le(m-1)R_1$, the last two terms
in \eqref{eq:first} have nonnegative sum. The function $g_m$ increases on
$0\le k<m$, and its average there is $(4m^2-9m-4)/3$. Thus
\[
 \T_1(N_{2m})-\T_1(A_{2m})
 \ge wR_0\frac{4m^2-9m-4}{3}>0\qquad(m\ge3).
\]

For the second position, put
\[
 U_z=\sum_{k=z}^m k,\qquad V_z=\sum_{k=z}^m k^2,\qquad
 W_z=\sum_{k=z}^m H_m(k)
\]
and define
\begin{equation}\label{eq:k2}
 K_m(l,t)=4V_lV_t-U_lW_t-U_tW_l
 -\frac{4m^2(m^2-1)}3-(m-l)(m-t).
\end{equation}
Let $e_0=w_0=1$ and $e_l=w_l-w_{l-1}\ge0$ for $l\ge1$.
The tail decomposition $w_k=\sum_{l\le k}e_l$ and polarization give
\begin{equation}\label{eq:comp2}
 \T_2(N_{2m})-\T_2(A_{2m})=\sum_{l,t=0}^m e_le_tK_m(l,t).
\end{equation}
Here $K_m(l,t)\ge0$ for all $0\le l\le t\le m$.
For an explicit certificate, substitute the four rows of
Table \ref{tab:k2}; each expansion of $36K_m(l,t)$ has only nonnegative
integer coefficients. The complete coefficient lists, totaling 105
monomials, are supplied in \texttt{kernel2.json}. Their exact verification
is explained in Section \ref{sec:verify}.
The four cases exhaust the domain according as $l\ge1$, or $l=0$ and
$t\ge2$, $t=0$, or $t=1$.
\begin{table}[ht]
\centering
\begin{tabular}{c c c c}
\toprule
Case & $l$ & $t$ & $m$\\\midrule
1 & $u+1$ & $u+1+b$ & $u+1+b+c$\\
2 & $0$ & $b+2$ & $b+2+c$\\
3 & $0$ & $0$ & $c$\\
4 & $0$ & $1$ & $c+1$\\\bottomrule
\end{tabular}
\caption{Second-position certificate substitutions. All free variables are
nonnegative integers.}\label{tab:k2}
\end{table}
In particular,
\[
 36K_m(0,0)=4m^6+24m^5+4m^4+48m^3+28m^2>0\quad(m\ge1).
\]
The $e_0^2$ term in \eqref{eq:comp2} proves strictness. Together with
\eqref{eq:apos}, this settles the first two positions.

\subsection{The third-position kernel}
For formal weights $w_0,\ldots,w_m$, set
$L_k=kw_k-\sum_{h<k}w_h$ and
$Q_k=2\sum_{h=0}^{\lfloor k/2\rfloor}h^2w_h$.
The quartic coefficient functionals extending those above are
\begin{align}
 c_4&=\sum_{k=1}^m\left[
 \binom{2k}{4}w_k+\binom{2k}{2}L_k+Q_k
 +2\sum_{h<k}h^2w_h+\ind_{2\mid k}\sum_{h<k}w_h\right],\label{eq:quartic}\\
 a_4&=\binom{2m}{4}w_m+\binom{2m}{2}L_m+Q_m.\notag
\end{align}
Indeed, the carry factor has $x^2$ coefficient $1$, $x^4$ coefficient
$\ind_{2\mid k}$, and no odd coefficients. Extract $x^4$ in
\eqref{eq:rec} and telescope.
For the formal tail vector $w_k=\ind_{k\ge z}$, denote the resulting
functionals by $p_j(z)$ and $a_j(z)$. Define the symmetric kernel
\begin{align}\label{eq:k3}
 2K^{(3)}_m(l,t)={}&2p_3(l)p_3(t)-p_2(l)p_4(t)-p_2(t)p_4(l)\notag\\
 &-2a_3(l)a_3(t)+a_2(l)a_4(t)+a_2(t)a_4(l).
\end{align}
This definition uses exactly the formal extension in
\eqref{eq:low}, \eqref{eq:cubic} and \eqref{eq:quartic}.

\begin{lemma}\label{lem:k3}
For integers $m\ge5$ and $0\le l\le t\le m$,
$K^{(3)}_m(l,t)\ge0$, and $K^{(3)}_m(0,0)>0$.
\end{lemma}
\begin{proof}
We give the piecewise polynomial formulas and the complete coverage argument
for an exact certificate. Put $q=\lfloor m/2\rfloor$, $\delta=m\bmod2$,
$\eta=z\bmod2$, and
\[
 v_r=\sum_{k=z}^m k^r,\qquad
 H_r=\begin{cases}\sum_{k=z}^q k^r,&z\le q,\\0,&z>q,\end{cases}
 \qquad E_z=\left(q-\frac{z-\eta}{2}\right)
 \left(q-\frac{z-\eta}{2}+1-\eta\right).
\]
Reversing the finite sums in \eqref{eq:quartic} gives
\begin{align}
 p_2(z)&=2v_2,\notag\\
 p_3(z)&=\tfrac43v_3-v_2+(2m+\tfrac53)v_1-(m^2+m)v_0,\notag\\
 p_4(z)&=\tfrac23v_4-2v_3+\tfrac{11}6v_2-\tfrac12v_1
       +z(2v_2-v_1)+2mv_2-2v_3\notag\\
       &\quad+2(m+1)H_2-4H_3+E_z,\label{eq:tails}\\
 a_2(z)&=2m^2-m+z,\notag\\
 a_3(z)&=(4m^3+2m)/3-2m(m-z),\notag\\
 a_4(z)&=\binom{2m}{4}+\binom{2m}{2}z+2H_2.\notag
\end{align}
For clarity, the four contributions to $p_4$ before expansion are
\[
 \sum_{k=z}^m\left[\binom{2k}{4}+z\binom{2k}{2}\right],\quad
 2\sum_{h=z}^q h^2(m-2h+1),\quad
 2\sum_{h=z}^m h^2(m-h),\quad E_z.
\]
The second sum is zero on its inactive branch, and
$E_z=\sum_{\substack{k\le m\\2\mid k}}(k-z)_+$.
The power sums are differences $F_r(m)-F_r(z-1)$, where
\begin{align*}
 F_0(y)&=y+1,& F_1(y)&=y(y+1)/2,& F_2(y)&=y(y+1)(2y+1)/6,\\
 F_3(y)&=y^2(y+1)^2/4,&
 F_4(y)&=y(y+1)(2y+1)(3y^2+3y-1)/30.
\end{align*}
Thus the formulas include $z=0$ and both sides of the halfway boundary.

Table \ref{tab:k3} gives chambers with $m=2(u+b+c)+\delta$.
For LL the inverse gaps $(u,b,c)$ are $(l,t-l,q-t)$; for HH they are
$(l-q-\delta,t-l,m-t)$. In the mixed case $l\le q<t$, use LH1 when
$t\ge q+l+\delta$, with inverse $(l,m-t,t-q-l-\delta)$, and otherwise
use LH2, with inverse $(t-q-\delta,q+l+\delta-t,q-l)$.
These formulas prove coverage.
\begin{table}[ht]
\centering\small
\begin{tabular}{c c c c c}
\toprule
Chamber & $l$ & $t$ & Extra condition & Active $H$ terms\\\midrule
LL & $u$ & $u+b$ & none & both\\
HH & $2u+b+c+\delta$ & $2u+2b+c+\delta$ & $u\ge1-\delta$ & neither\\
LH1 & $u$ & $2u+b+2c+\delta$ & $u+c+\delta\ge1$ & $l$ only\\
LH2 & $b+u$ & $b+2u+c+\delta$ & $u\ge1-\delta$ & $l$ only\\\bottomrule
\end{tabular}
\caption{Third-position chambers, with $u,b,c\ge0$ and
$\delta\in\{0,1\}$.}\label{tab:k3}
\end{table}

For $\delta=0$, split LH1 into $u\ge1$ and $u=0,c\ge1$.
Shift all required lower bounds, then split each free gap by parity.
The resulting 68 roots have gaps $g_i=g_i^0+2X_i$, with $X_i\ge0$;
the fixed first coordinate in the second LH1 subcase is zero.
The remaining restriction $m\ge5$ is
\begin{equation}\label{eq:threshold}
 X_0+X_1+X_2\ge
 \max\left(0,\left\lceil\frac{3-\delta-\sum_i g_i^0}{2}\right\rceil\right),
\end{equation}
whose right side is at most $2$.
The supplementary file \texttt{kernel3.json} lists 144 leaves. A leaf
specifies its root, an offset $o_i$, and whether each axis is fixed;
on a free axis $X_i=o_i+Y_i$, $Y_i\ge0$, and on a fixed axis $X_i=o_i$.
For every leaf the supplied list is the complete expansion
\begin{equation}\label{eq:cert}
 1440K^{(3)}_m(l,t)=\sum_{\boldsymbol a}
 C_{\boldsymbol a}Y_0^{a_0}Y_1^{a_1}Y_2^{a_2},\qquad
 C_{\boldsymbol a}\in\mathbb Z_{\ge0}.
\end{equation}
There are 12,325 listed monomials; omitted coefficients are zero.
Section \ref{sec:verify} gives a direct exact identity test for every row.
Every offset is at most $2$. Clipping $X_i$ to $\min(X_i,4)$ preserves
membership in each leaf and preserves \eqref{eq:threshold}.
The finite check on these clipped vectors covers all 68 roots and all
8,059 valid vectors. Each whole leaf satisfies the domain, parity and
halfway inequalities, as checked at its origin and on its free directions.
Thus \eqref{eq:cert} proves nonnegativity on the entire infinite domain.
Finally, the two expansions of $1440K^{(3)}_m(0,0)$ for $m=2z+5$ and
$m=2z+6$ have strictly positive coefficients; they are displayed in
Appendix \ref{app:strict}. This proves strictness.
\end{proof}

\begin{proof}[Completion of Theorem \ref{thm:main}(2)]
The same tail decomposition used in \eqref{eq:comp2} yields
\[
 \T_3(N_{2m})-\T_3(A_{2m})
 =\sum_{l,t=0}^m e_le_tK^{(3)}_m(l,t)>0\qquad(m\ge5),
\]
by Lemma \ref{lem:k3} and $e_0=1$.
It remains to prove $\T_3(A_{2m})>0$ separately. Let
$L=mw-R_0\ge0$, $C_i=\binom{2m}{i}$, and
$q_2=2\sum_{h=0}^{\lfloor m/2\rfloor}h^2$.
Monotonicity gives $Q_m\le q_2w$. Since
\[
 a_2=C_2w+L,\quad a_3=C_3w+2mL,\quad
 a_4=C_4w+C_2L+Q_m,
\]
we have
\begin{align}\label{eq:thirdA}
 \T_3(A_{2m})&\ge a(m)w^2+b(m)wL+(2m^2+m)L^2,\\
 a(m)&=C_3^2-C_2(C_4+q_2),\qquad
 b(m)=4mC_3-C_2^2-C_4-q_2.\notag
\end{align}
The expansions of $9a(m)$ and $3b(m)$ at $m=2z+4$ and $m=2z+5$
have strictly positive coefficients (Appendix \ref{app:strict}).
Hence \eqref{eq:thirdA} is positive for $m\ge4$.
For the remaining comparisons the exact triples are
\[
\begin{array}{c|c|c|c}
 n&(c_2,c_3,c_4)&(a_2,a_3,a_4)&\T_3(N_n)>\T_3(A_n)\\\hline
 6&(54,80,100)&(32,52,62)&1000>720\\
 8&(182,406,738)&(122,304,578)&30520>21900.
\end{array}
\]
This settles every even $n\ge6$. Both $A_{2m}$ and $N_{2m}$ are
palindromic of degree $2m+2d(m)=d(2m)$, so the comparisons reflect.
The identity $N_{2m+1}=N_{2m}$ proves the claimed odd-index statements.
\end{proof}

\section{Exact verification and limitations}\label{sec:verify}
The supplementary directory contains the complete polynomial coefficients
and a stand-alone Python 3 checker, \texttt{verify.py}, which imports no
symbolic algebra package. Its checks are polynomial identities, rather
than sampled positivity assertions. On each fixed parity and halfway chamber,
\eqref{eq:tails} gives total degrees at most $3,4,5$ for $p_2,p_3,p_4$
and $2,3,4$ for $a_2,a_3,a_4$. Thus $K^{(3)}$ has total degree at most $8$
after every affine leaf substitution; $K$ in \eqref{eq:k2} has degree at most
$6$. The checker verifies these bounds on the supplied expansions and
compares them to the literal defining sums on every nonnegative integer
simplex point of total degree at most $8$ or $6$, respectively.

For completeness, a polynomial of total degree at most $D$ is determined
by its values at $\boldsymbol v\in\mathbb Z_{\ge0}^r$ with
$|\boldsymbol v|\le D$. Expand it in the basis
$\prod_i\binom{X_i}{a_i}$, $|\boldsymbol a|\le D$.
The coefficients are mixed forward differences at the origin; each uses
only values at $\boldsymbol v\le\boldsymbol a$ and hence only this simplex.
Agreement there is therefore an identity. This is the lower-set Newton
interpolation principle; see \cite{Neidinger} for its general formulation.
The present finite-difference argument supplies the precise special case used.

The resulting checks comprise 144 third-position identities at 12,720
exact simplex evaluations, four second-position identities at 126 evaluations,
and six strictness identities at 42 univariate evaluations. The separate
coverage check is essential: polynomial identity verification alone would
not show that the leaves exhaust the allowed parameters.
The original expansions were generated with SymPy \cite{SymPy} and also
checked by a separate rational polynomial implementation.
Related exact verification of algebraic positivity occurs in
Section 7 of \cite{ZengChollet}, where symbolic integer and rational
calculations check sum-of-squares and positive-semidefinite Gram identities
for a restricted matrix inequality. Our certificates use nonnegative monomial
expansions on the unbounded chambers above.

The rational test \eqref{eq:test} can be evaluated without computing either
$b(n)$ or $N_n$. For example, exact evaluations at $n=2^r$ give the following
initial consecutive ranges. These are applications of the proved sufficient
criterion, not coefficient enumerations at those parameters.
\[
\begin{array}{c|rrrrrrrr}
 r&8&16&32&64&128&256&512&1024\\\hline
 \text{positions }1\le j\le&1&2&3&4&7&10&15&22
\end{array}
\]
A second supplementary script, \texttt{criterion.py}, reproduces this table
using exact fractions. Finite calibration during the derivation checked the
mean bound for $2\le n\le4096$ and coefficient envelopes through degree
$13$; it does not replace any infinite argument above.
The range between the proven boundary positions remains uncontrolled.
In particular, neither full unimodality nor \eqref{eq:original} follows from
these results. The sufficient test may be conservative, and no optimal
threshold $n_\varepsilon$ is claimed.

\paragraph{Computational and AI assistance.}
OpenAI Codex assisted with mathematical derivations, symbolic certificate
construction, exact verification code, literature checking, and preparation
of this manuscript. A separate model context reconstructed the submitted
arguments and checked the certificates from their defining sums using
the interpolation method above. Such a pass is not external peer review.
All certificate sign decisions described here use integers or rational
numbers. No Lean or other proof-assistant formalization is claimed.

\appendix
\section{Strictness polynomials}\label{app:strict}
In each row below the displayed list gives the coefficients in ascending
powers of $z\ge0$; line breaks do not start a new list. All are positive.
The supplementary \texttt{strictness.json} contains the same arrays.
\begin{align*}
1440K^{(3)}_{2z+6}(0,0):\quad&[46638720,139543488,175898688,123122688,\\
 &52593728,14069952,2302592,210432,8192],\\
1440K^{(3)}_{2z+5}(0,0):\quad&[8985600,34793280,55319456,48077056,\\
 &25219648,8213440,1623424,177664,8192],\\
9a(2z+4):\quad&[8064,28632,41224,31014,12916,2832,256],\\
9a(2z+5):\quad&[40500,105300,113725,65304,21028,3600,256],\\
3b(2z+4):\quad&[96,406,463,206,32],\\
3b(2z+5):\quad&[465,1056,823,270,32].
\end{align*}
The first two polynomial identities have degree $8$, the next two degree
$6$, and the last two degree $4$. Their exact univariate verification uses
$9$, $7$, and $5$ integer arguments per row, respectively. The degree
bounds follow from \eqref{eq:tails} and \eqref{eq:thirdA}.


\raggedright
\bibliographystyle{plainnat}
\bibliography{references}

\end{document}
